\chapter{Methodology}
\label{chap:methodology}

\section{Research design}
Name and justify the research design in relation to the research questions. Explain
why it is a better fit than plausible alternatives, and distinguish the overall
methodology from individual data-collection methods.

\section{Study workflow}
\begin{figure}[htbp]
  \centering
  \begin{tikzpicture}[
  node distance=9mm,
  every node/.style={draw,rounded corners,align=center,minimum width=3.6cm,minimum height=8mm},
  every path/.style={draw,-{Latex[length=2mm]}}
]
  \node (questions) {Research questions};
  \node (evidence) [below=of questions] {Evidence collection};
  \node (analysis) [below=of evidence] {Analysis};
  \node (validation) [below=of analysis] {Validation and interpretation};
  \path (questions) -- (evidence);
  \path (evidence) -- (analysis);
  \path (analysis) -- (validation);
  \path (validation.east) -- ++(1.0,0) |- (questions.east);
\end{tikzpicture}


  \caption{Example research workflow drawn directly in \LaTeX.}
  \label{fig:research-workflow}
\end{figure}

\CorrectionAdd{E2-01}{a}{As illustrated in \Cref{fig:research-workflow}, the research
questions determine the evidence required, the analysis converts that evidence into
findings, and validation checks whether the resulting interpretation is sufficiently
supported. The arrows represent an iterative process rather than an irreversible
pipeline.}

\section{Data collection or evidence generation}
Describe who or what supplied the evidence, how items were selected, the procedure,
instrument quality and any departures from the approved protocol. Replace this example
with enough detail for a knowledgeable reader to assess validity and reproducibility.

\section{Analysis}
Define every measure and connect it to the question it answers. For example, an
illustrative normalised change can be expressed as
\begin{equation}
  \Delta_{i}=\frac{x_{i,\mathrm{post}}-x_{i,\mathrm{pre}}}{
  \max(x_{i,\mathrm{pre}},\epsilon)},
  \label{eq:normalised-change}
\end{equation}
where $x_{i,\mathrm{pre}}$ and $x_{i,\mathrm{post}}$ are the verified measurements for
case $i$, and $\epsilon>0$ prevents division by zero. Equation~\eqref{eq:normalised-change}
is an example only; students must use an analysis justified for their own evidence.

\section{Ethics, integrity and data management}
State the relevant approval identifier, consent arrangements, privacy safeguards,
storage and retention controls. Never copy an example approval number. Follow the
current instructions from your institution, faculty or school and research ethics
office.

